\chapter*{Testsysteme}
\addcontentsline{toc}{chapter}{Testsysteme}

Zur Evaluation der Lastkriterien wurden zwei Referenzsysteme verwendet, die in Tabelle~\ref{tab:hardware} aufgeführt sind.
Die Systeme werden im Dokument der Einfachheit halber abkürzend als \textit{Desktop} und \textit{Laptop} bezeichnet.\\

Alle Benchmarks wurden unter Windows 11 Pro (x64) durchgeführt, wobei die \textit{Release-Builds} der jeweils getesteten Softwarekomponenten zum Einsatz kamen.\\

\begin{table}[!htbp]
    \centering
    \begin{tabularx}{\textwidth}{l Y Y}
    \toprule
    \textbf{} & \textbf{Desktop} & \textbf{Laptop} \\
    \midrule
    \textbf{CPU}  & AMD Ryzen 9 9950X3D (16 Kerne, 32 Threads, 4\,300\,MHz) & Intel Core i7-8750H (6 Kerne, 12 Threads) \\
    \textbf{RAM}  & 64\,GB DDR5       & 16\,GB DDR4 \\
    \textbf{GPU}  & NVIDIA GeForce RTX 5090 & NVIDIA GeForce GTX 1070 Max-Q \\
    \bottomrule
    \end{tabularx}
    \caption{Spezifikationen der verwendeten Testsysteme.}
    \label{tab:hardware}
\end{table}

\subsection*{Cache-Hierarchie}
Tabelle~\ref{tab:cache_hierarchy} zeigt die jeweilige Cache-Hierarchie beider Systeme.

\begin{table}[!htbp]
    \centering
    \begin{tabularx}{\textwidth}{l Y Y}
    \toprule
    \textbf{} & \textbf{Desktop} & \textbf{Laptop} \\
    \midrule
    L1 Data        & 48\,KiB $\times$ 16 & 32\,KiB $\times$ 6 \\
    L1 Instruction & 32\,KiB $\times$ 16 & 32\,KiB $\times$ 6 \\
    L2 Unified     & 1\,024\,KiB $\times$ 16 & 256\,KiB $\times$ 6 \\
    L3 Unified     & 128\,MiB & 9\,MiB  \\
    \bottomrule
    \end{tabularx}
    \caption{Cache-Hierarchie der verwendeten Testsysteme.}
    \label{tab:cache_hierarchy}
\end{table}

